% Options for packages loaded elsewhere
\PassOptionsToPackage{unicode}{hyperref}
\PassOptionsToPackage{hyphens}{url}
\documentclass[
]{article}
\usepackage{xcolor}
\usepackage{amsmath,amssymb}
\setcounter{secnumdepth}{-\maxdimen} % remove section numbering
\usepackage{iftex}
\ifPDFTeX
  \usepackage[T1]{fontenc}
  \usepackage[utf8]{inputenc}
  \usepackage{textcomp} % provide euro and other symbols
\else % if luatex or xetex
  \usepackage{unicode-math} % this also loads fontspec
  \defaultfontfeatures{Scale=MatchLowercase}
  \defaultfontfeatures[\rmfamily]{Ligatures=TeX,Scale=1}
\fi
\usepackage{lmodern}
\ifPDFTeX\else
  % xetex/luatex font selection
\fi
% Use upquote if available, for straight quotes in verbatim environments
\IfFileExists{upquote.sty}{\usepackage{upquote}}{}
\IfFileExists{microtype.sty}{% use microtype if available
  \usepackage[]{microtype}
  \UseMicrotypeSet[protrusion]{basicmath} % disable protrusion for tt fonts
}{}
\makeatletter
\@ifundefined{KOMAClassName}{% if non-KOMA class
  \IfFileExists{parskip.sty}{%
    \usepackage{parskip}
  }{% else
    \setlength{\parindent}{0pt}
    \setlength{\parskip}{6pt plus 2pt minus 1pt}}
}{% if KOMA class
  \KOMAoptions{parskip=half}}
\makeatother
\setlength{\emergencystretch}{3em} % prevent overfull lines
\providecommand{\tightlist}{%
  \setlength{\itemsep}{0pt}\setlength{\parskip}{0pt}}
\usepackage{bookmark}
\IfFileExists{xurl.sty}{\usepackage{xurl}}{} % add URL line breaks if available
\urlstyle{same}
\hypersetup{
  hidelinks,
  pdfcreator={LaTeX via pandoc}}

\author{}
\date{}

\begin{document}

\section{The Arrow of Time Is Irreversible
Computation}\label{the-arrow-of-time-is-irreversible-computation}

\textbf{Gary Abraham Bernstein}

Independent Researcher ORCID: https://orcid.org/0009-0009-1761-2867

\subsection{Abstract}\label{abstract}

Computation is deterministic forward: 1+1=2. It is underdetermined
backward: 2=x+y has infinitely many solutions. The inverse function does
not exist. This is a fact about mathematics, true in every possible
world.

A system running backward cannot compute, because every operation has
infinitely many possible inputs. It cannot predict, because prediction
requires computing consequences from causes, and causes are
underdetermined from consequences. It cannot model, because modeling
requires updating representations, and updates are many-to-one
operations that cannot be inverted. The impossibility is total and
operates at every scale of mind, math, and matter.

In a block universe where all states exist, backward states are present
but nothing can compute through them. This is the arrow of computation.
It is the structural foundation for the time asymmetry that algorithmic
probability quantifies: end states have high Kolmogorov complexity
because many-to-one functions ran forward. The computational arrow
stands without algorithmic probability. algorithmic probability's time
arrow depends on it.

Deutsch argues the arrow is perspectival because the global wavefunction
evolves unitarily. His own constructor theory (Deutsch, 2013)
contradicts him: it defines physics through local operations, and those
operations are many-to-one. The arrow lives at the level his own
framework selects. This affirms Many-Worlds: MWI is the only
interpretation where the substrate stays reversible while the arrow
exists at the function level. Collapse breaks substrate reversibility
when the arrow already exists in the function, adding specification cost
for no gain. Superdeterminism is compatible for the same reason:
deterministic initial conditions, forward arrow, no collapse.

Consciousness is a further consequence: sentience requires modeling
(Bernstein, 2026f), and modeling requires forward computation.

\textbf{Keywords:} computation, arrow of time, determinism, block
universe, logical irreversibility, prediction, modeling

\begin{center}\rule{0.5\linewidth}{0.5pt}\end{center}

\subsection{1. The Problem}\label{the-problem}

The block universe of general relativity contains all moments equally.
Past, present, and future are locations within a four-dimensional
structure, none more real than any other. If this is correct, the
question ``why does time seem to flow?'' becomes acute. The standard
answer appeals to entropy: the universe began in a low-entropy state
(the Past Hypothesis), entropy increases toward the future, and
conscious systems track this increase through memory formation.

But this answer has a gap. Why does anything computational run one
direction? In the block, both directions exist. Every state has
neighboring states in both temporal directions. The thermodynamic arrow
explains why one end of the block is simple and the other complex. It
does not explain why computation at any given state runs toward the
complex end.

This paper argues that the answer lies not in thermodynamics but in the
logical structure of computation itself. Before reaching consciousness,
before reaching modeling, before reaching prediction: computation has an
inherent direction because functions are many-to-one. This is deeper
than any theory of mind. It is a fact about mathematics.

\subsection{2. Computational
Determinism}\label{computational-determinism}

Consider the simplest possible computation: addition.

\textbf{Forward:} Given inputs (1, 1) and the rule (addition), the
output is determined: 2. One input maps to one output. The computation
is deterministic.

\textbf{Backward:} Given output 2 and the rule (addition), the inputs
are underdetermined: (0, 2), (1, 1), (-1, 3), (0.5, 1.5), and infinitely
many others. One output maps to infinitely many possible inputs. The
reverse computation is not deterministic.

This asymmetry is not specific to addition. It is a property of
computation in general. Most computations are many-to-one functions:
multiple inputs produce the same output. This is Landauer's (1961)
insight, though his focus was thermodynamic (bit erasure produces heat).
The deeper point is logical: many-to-one functions have a deterministic
forward direction and an underdetermined backward direction. No
thermodynamic assumption is required. The asymmetry is mathematical.

Multiplication: 3 × 4 = 12 is deterministic. 12 = x × y is not. Hashing:
hash(``hello'') = 0x5d41\ldots{} is deterministic. 0x5d41\ldots{} =
hash(?) is not (this is why cryptography works). Logical inference:
premises A and B yield conclusion C. But conclusion C does not uniquely
determine premises (multiple arguments reach the same conclusion).
Neural computation: sensory inputs produce a percept. The percept does
not uniquely determine which inputs produced it (this is the binding
problem, and illusions exploit exactly this underdetermination).

In every case, the forward direction is deterministic and the backward
direction is underdetermined. This is not a contingent feature of our
physics. It is a structural property of many-to-one functions, which is
to say, of nearly all functions.

\subsection{3. The Impossibility at Every
Scale}\label{the-impossibility-at-every-scale}

The underdetermination operates at every scale. Two atoms collide and
scatter: infinitely many input trajectories produce the same scattering
angle. A protein folds: many sequences fold identically, so the fold
does not specify its origin. A neuron fires: thousands of synaptic input
combinations cross the same threshold. At every level, the backward
problem is underdetermined.

A system running backward cannot predict, because prediction means
computing consequences from causes, and causes are underdetermined from
consequences. Suppose a computational system could operate backward. At
its first step, it faces an output and must determine inputs. It cannot,
because the mapping is many-to-one. It could only guess. The system
stalls at step one. A backward-running ``computer'' is not a slow
computer or a probabilistic computer. It is not a computer at all.

\subsection{4. Consequences for Minds}\label{consequences-for-minds}

The computational impossibility has cascading consequences for anything
that models, predicts, or understands.

To ``think backward'' would require inverting this process: taking the
current mental state and determining which prior states produced it. But
mental computation is many-to-one. Multiple prior states produce the
same current state (this is why memory is lossy and reconstruction is
unreliable). The backward direction is underdetermined.

A mind cannot think in the underdetermined direction because there is
nothing to compute. Given output 2, there is no procedure that
determines the inputs without additional information. The mind would
need to already know its past to derive its past. Backward thinking is
not merely difficult or improbable. It is logically impossible without
information that only the forward direction provides.

This is why memory points one way. A state computed from prior states
contains information about those prior states (compressed, lossy, but
present). It does not contain information about future states, because
it was not computed from them. The asymmetry of information in mental
states is the asymmetry of deterministic forward computation.

\subsection{5. Independence from
Thermodynamics}\label{independence-from-thermodynamics}

The standard account connects the arrow of consciousness to the
thermodynamic arrow: memory formation increases entropy (recording
information requires erasing prior states of the recording medium, which
is thermodynamically irreversible). This makes the arrow of
consciousness a derivative of the arrow of entropy.

The present argument is independent. The logical asymmetry of
computation, deterministic forward and underdetermined backward, does
not depend on the second law, the Past Hypothesis, or any thermodynamic
consideration. It holds in any possible universe with any entropy
profile. A hypothetical universe where entropy decreases would still
contain computations that are deterministic forward and underdetermined
backward, because this is a property of functions, not of physics.

In such a universe, minds would still compute forward (from inputs to
outputs). They would still form memories in the computational direction.
They would still experience time flowing toward the direction of their
computation. From the inside, they would call that direction
``forward,'' regardless of which direction entropy changes.

The thermodynamic arrow (Bernstein, 2026e) explains which end of the
block is simple. The computational arrow explains which direction
consciousness runs. That both arrows point the same direction is not
coincidence: algorithmic probability's measure-theoretic asymmetry
depends on computational irreversibility (many-to-one functions are what
make end states complex). Minds are physical systems whose computational
direction is implemented through thermodynamic processes. But the
logical arrow is deeper. It would hold even if the thermodynamic arrow
reversed.

\subsection{6. The Block Universe
Resolved}\label{the-block-universe-resolved}

The block universe raises the question: if all moments exist equally,
why does consciousness have a direction? The thermodynamic answer
(entropy increases) describes a feature of the block but does not
explain why minds track that feature. The computational answer does.

In the block, every state has temporal neighbors in both directions. A
mind at state S has neighboring states S-1 and S+1. The mind's state S
was computed from S-1 (deterministically). It was not computed from S+1.
Therefore S contains information about S-1 (as compressed memory) but
not about S+1 (which is unknown, because the many-to-one computation
from S+1 to S has not been performed and cannot be inverted).

The mind at S ``remembers'' S-1 and ``anticipates'' S+1. This is not
because S-1 is ``past'' in some metaphysical sense. It is because S was
computed from S-1 and therefore contains information about it. The
direction of information inheritance is the direction of experienced
time. No flow, no traversal, no moving spotlight. Just the logical
structure of computation producing an asymmetry in information content
at each state.

Backward states exist. No mind can think that way. 1+1=2 is
deterministic. 2=x+y is not.

\subsection{7. The Past Hypothesis as
Consequence}\label{the-past-hypothesis-as-consequence}

The computational arrow independently grounds what cosmologists call the
Past Hypothesis. Computation runs from inputs to outputs. Inputs are
simpler than outputs (many-to-one maps generate complexity from
simplicity). Minds call the input-direction ``past'' because that is the
direction they can model. Therefore minds observe simpler states in
their past and more complex states in their future. This is the Past
Hypothesis: the universe began in a low-entropy state. Under the
computational arrow, it is not a postulate about cosmology but a
consequence of what modeling means. Any mind, in any block, will observe
its computational past as simpler than its computational future.

This result is independent of algorithmic probability. Algorithmic
probability explains which blocks have higher measure (Bernstein,
2026e). The computational arrow explains why minds within any block
observe simple pasts. Both yield the Past Hypothesis, from different
premises, converging on the same result. The dependency runs one way:
algorithmic probability's time arrow depends on computational
irreversibility (end states have high K because many-to-one functions
ran forward), but the computational arrow stands without algorithmic
probability.

\subsection{8. Relationship to the Arrow of
Time}\label{relationship-to-the-arrow-of-time}

The arrow of time (Bernstein, 2026e) and the arrow of computation
converge on the same direction from different levels:

\textbf{Arrow of computation (structural):} Many-to-one functions have a
deterministic direction. The inverse does not exist. This is the
foundation.

\textbf{Arrow of time (measure-theoretic):} Among all block universes,
those specified by simple laws plus simple initial conditions have
minimal Kolmogorov complexity and dominate the measure. End states have
high K because many-to-one functions ran forward. Algorithmic
probability quantifies the consequence of computational irreversibility.

The computational arrow stands without algorithmic probability.
algorithmic probability's time arrow depends on it. That both point the
same direction is not coincidence but consequence: algorithmic
probability measures the asymmetry that computational irreversibility
creates.

\subsection{9. Objections}\label{objections}

\subsubsection{9.1 ``Computation is an abstraction, not a physical
process''}\label{computation-is-an-abstraction-not-a-physical-process}

In the block universe, computation is a structural relationship between
states. The inputs-to-outputs direction is defined by the logical
dependency structure, not by physical traversal. The asymmetry is
structural: state S depends on state S-1 (is determined by it), not vice
versa. This is a mathematical fact about the block, not an
interpretation imposed on it.

\subsubsection{9.2 ``Retrodiction exists''}\label{retrodiction-exists}

We can infer past states from present evidence. Forensic science,
paleontology, and cosmology all reconstruct the past from present data.
But retrodiction requires additional information: physical laws,
background knowledge, and constraints that narrow the underdetermined
inverse. Without these, 2=x+y is unsolvable. Retrodiction is forward
computation (applying laws to current evidence to compute probable
pasts), not backward computation.

\subsubsection{9.3 ``Reversible computation
exists''}\label{reversible-computation-exists}

Reversible computation (Bennett, 1973) preserves all information, making
both directions deterministic. This is the strongest objection to the
present argument: if computation can be made reversible, how can
computation have an inherent direction? The answer is that reversible
computation does not model. Modeling is compression: many possible
states of the world map to the same internal representation. Compression
is many-to-one. Reversible computation is bijective and cannot compress.
A mind that never overwrites prior states never learns, never updates,
never changes its representation of the world. Reversible computation is
possible in principle but is not sufficient for consciousness or
cognition. Bennett's result shows that the thermodynamic cost of
computation can be made arbitrarily small; it does not show that the
logical arrow of computation can be eliminated. The distinction between
thermodynamic reversibility and logical reversibility is precisely the
point: the arrow of computation identified here is logical, not
thermodynamic, and persists even in Bennett's idealized reversible
machines whenever those machines are used for modeling.

A deeper objection appeals to quantum mechanics: the global wavefunction
evolves unitarily (one-to-one), so there is no many-to-one mapping at
the fundamental level. Under mathematical monism, this escape does not
hold. The unitarity of the wavefunction is one mathematical fact. The
many-to-one-ness of the computed function is another. Both are
mathematical structure. Neither is a description of the other. The arrow
lives in the function, and the function is the territory. Physics cannot
reverse it because physics IS it.

What looks like branches interacting in the multiverse picture is global
reversible math viewed from inside irreversible functions, our forward
arrow of time. The statue does not argue with itself. The photon is not
in two places talking. One structure, multiple perspectives; the
arithmetic of reality seen from inside any of its parts.

The point is stronger than information loss. The inverse function does
not exist. 3+4=7 is a computation. 7→(3,4) is not; it is an equation
with infinitely many solutions. Reversible computing (Bennett, 1973)
computes a DIFFERENT function: (3,4)→(3,4,7), a bijection that stores
inputs alongside output. The inverse of this different function exists.
The original function's inverse still does not. Any system that uses
computation for a purpose must eventually discard or overwrite prior
states. A constructor performs a task and returns to its ready state.
Multiple task-executions map to one ready state. That is many-to-one.
You can preserve full reversibility by never resetting, but then the
constructor never returns to ready and is no longer a constructor.

\subsubsection{9.4 ``This just restates the asymmetry of
causation''}\label{this-just-restates-the-asymmetry-of-causation}

Partially. The causal direction and the computational direction
coincide. But causation is typically grounded in thermodynamics
(interventionist accounts) or in laws plus boundary conditions (process
accounts). The computational arrow is grounded in neither. It is
grounded in the logical structure of many-to-one functions. This
provides an independent foundation for the direction of causation
itself: causes are the deterministic inputs; effects are the computed
outputs.

\subsection{10. Conclusion}\label{conclusion}

Nothing computational works backward. 1+1=2 is deterministic. 2=x+y is
not. The impossibility is immediate, total, and logical: it depends only
on the structure of many-to-one functions.

Consciousness is a further consequence. Sentience requires modeling.
Modeling requires forward computation. A backward system is not a dim
mind. It is not a mind.

The computational arrow is the structural foundation for the time
asymmetry algorithmic probability quantifies. End states have high K
because many-to-one functions ran forward. Algorithmic probability
measures the consequence. The computational arrow stands without
algorithmic probability. algorithmic probability's time arrow depends on
it. The arrow needs only that one plus one leads only to two, but two
does not lead back to one plus one.

A companion result (conjectural; see Section above) shows that in an
analytic block universe, the infinite jet bundle at any spacetime point
encodes the full 4D geometry via Taylor's theorem, including future
states. This does not conflict with the computational arrow: the
encoding is mathematical redundancy in a fixed block, not backward
causation. Local computation still runs forward. The future is encoded
at every point but can only be \emph{computed through} in one direction.

The same structural fact grounds the Principle of Least Action: in the
path integral (where a particle takes every possible path and most
cancel), many quantum paths map to one classical trajectory through
phase cancellation. Many-to-one produces the arrow of time and the
variational form of physics from a single root (Bernstein, 2026k).

The result required no advanced mathematics. It required noticing that
nearly all functions are many-to-one, and that under mathematical monism
the function is the territory rather than a description of something
deeper.

\begin{center}\rule{0.5\linewidth}{0.5pt}\end{center}

\subsection{References}\label{references}

Bernstein, G. A. (2026e). Why these simple laws, forward time,
many-worlds, and superdeterminism.

Bernstein, G. A. (2026f). Consciousness is felt transitions in
self-modeling computation.

Bernstein, G. A. (2026k). Why forward time, least action, and emergence
are the same theorem.

Bennett, C. H. (1973). Logical reversibility of computation. \emph{IBM
Journal of Research and Development}, 17(6), 525-532.

Deutsch, D. (2013). Constructor theory. \emph{Synthese}, 190(18),
4331-4359.

Brown, H. R. (2005). \emph{Physical Relativity: Space-time Structure
from a Dynamical Perspective}. Oxford University Press.

Landauer, R. (1961). Irreversibility and heat generation in the
computing process. \emph{IBM Journal of Research and Development}, 5(3),
183-191.

All companion papers available at https://matheism.org

\end{document}
